% !TEX TS-program = lualatex

\DocumentMetadata{
  lang = en,
  pdfversion = 2.0,
  pdfstandard = ua-2,
  tagging = on,
  tagging-setup = {
    math/setup = {mathml-SE},
    table/header-rows = 1
  }
}

\documentclass[12pt]{article}

\usepackage{amsmath}
\usepackage{geometry}
\geometry{
  letterpaper,
  textwidth=6.5in,
  textheight=9in,
  top=0.85in,
  bottom=1in,
  left=0.75in,
  right=0.75in,
}
\usepackage{hyperref}

\newtheorem{theorem}{Theorem}
\newtheorem{acknowledgement}[theorem]{Acknowledgement}
\newtheorem{algorithm}[theorem]{Algorithm}
\newtheorem{axiom}[theorem]{Axiom}
\newtheorem{case}[theorem]{Case}
\newtheorem{claim}[theorem]{Claim}
\newtheorem{conclusion}[theorem]{Conclusion}
\newtheorem{condition}[theorem]{Condition}
\newtheorem{conjecture}[theorem]{Conjecture}
\newtheorem{corollary}[theorem]{Corollary}
\newtheorem{criterion}[theorem]{Criterion}
\newtheorem{definition}[theorem]{Definition}
\newtheorem{example}[theorem]{Example}
\newtheorem{exercise}[theorem]{Exercise}
\newtheorem{lemma}[theorem]{Lemma}
\newtheorem{notation}[theorem]{Notation}
\newtheorem{problem}[theorem]{Problem}
\newtheorem{proposition}[theorem]{Proposition}
\newtheorem{remark}[theorem]{Remark}
\newtheorem{solution}[theorem]{Solution}
\newtheorem{summary}[theorem]{Summary}
\newenvironment{proof}[1][Proof]{\textbf{#1.} }{\ \rule{0.5em}{0.5em}}

\tagpdfsetup{
  math/alt/use,
  role/new-tag=section/H1,
}

\hypersetup{
  pdftitle={Assignment 3},
  pdfauthor={Blankenau},
  pdfsubject={Fall, 2027},
  hidelinks,
}

\begin{document}
\begin{center}
{\LARGE Assignment 3}
\end{center}

\noindent Blankenau

\noindent Economics 905, Fall, 2027

\medskip
\noindent \textit{Instructions.} You may work in groups and may compare
answers prior to submission. Joint submissions are allowed and encouraged.
All work should be legible and concise. A due date will be announced in
class.

\begin{enumerate}
\item Parts a-g follow your notes but try to work this without reference to
that for practice. Consider a one-period economy populated by a
representative agent and a representative firm. The agent is endowed with
one unit of time which is allocated between working, $n_{s},$ and leisure $%
\ell $. The agent makes these choices to maximize utility given by 
\begin{equation*}
u\left( c,\ell \right) =\frac{c^{1-\gamma }-1}{1-\gamma }+\frac{\ell
^{1-\gamma }-1}{1-\gamma }
\end{equation*}%
where $c$ is consumption and $\gamma >0$ measures the degree of curvature in
the utility function with respect to consumption and leisure. For each unit
of time worked, the agent is paid a wage $\omega $. The agent is also
endowed with $k_{s}$ units of capital which she rents inelastically in the
capital market. Each unit of capital rented pays $r$ and then fully
depreciates. The firm combines capital, $k_{d},$ and labor, $n_{d},$ to
produce a final good according to 
\begin{equation}
y\left( k,n\right) =zk_{d}^{\alpha }n_{d}^{1-\alpha },\text{ }0<\alpha <1
\label{out}
\end{equation}%
All markets are competitive.

\begin{enumerate}
\item Find the first order condition(s) for the agent's problem.

\item Find the first order conditions for the firm's problem.

\item Find the market clearing conditions.

\item Define an equilibrium.

\item Arrive at one equation and one unknown that could be used to solve for 
$\ell .$

\item Discuss how to find the other items of interest.

\item State and solve the social planner's problem. You will not be able to
get an explicit solution so arrive at one equation and 1 unknown. \ Will $%
\ell $ be the same as in the competitive equilibrium?

\item Suppose instead that there are two goods and preferences are given by 
\begin{equation*}
u\left( c,\ell \right) =\frac{c_{1}^{1-\gamma }-1}{1-\gamma }+\frac{%
c_{2}^{1-\gamma }-1}{1-\gamma }+\frac{\ell ^{1-\gamma }-1}{1-\gamma }
\end{equation*}%
where $c_{1}$ and $c_{2}$ are consumption of good 1 and good two. There is a
representative competitive firm producing each of the two goods according to 
\begin{eqnarray*}
y_{1} &=&z_{1}k_{1}^{\alpha }n_{1}^{1-\alpha } \\
y_{2} &=&z_{2}k_{2}^{\alpha }n_{2}^{1-\alpha }
\end{eqnarray*}%
where $n_{j}$ are $k_{j}$ are capital and labor hired by the firm producing
good $j$. Note that in equilibrium $k_{1}+k_{2}=k_{s}$ and $%
n_{1}+n_{2}=n_{s}.$ Write down the social planner's problem.

\item Find a system of three equations and three unknowns that could be used
to solve for the equilibrium values of $n_{1},n_{2},$ and $k_{1}.$

\item For the case where $\gamma =1$ find $n_{1},$ $n_{2},$ $\ell ,$ $k_{1}$
and $k_{2}.$
\end{enumerate}

\item Consider the following representative agent model. The consumer has
preferences given by 
\begin{equation*}
u\left( c,\ell \right) =\ln c+\ell
\end{equation*}%
where $c$ and $\ell $ are consumption and leisure$.$ The consumer is endowed
with 1 unit of time and $k_{0}$ units of capital. The representative firm
has a technology for producing consumption goods given by $y=zk^{\gamma
}n^{1-\gamma }$ where $z$ is total factor productivity, $k$ is the capital
input, $n$ is the labor input and $0<\gamma <1.$ Let $\omega $ denote the
wage rate and $r$ the rental rate on capital.

\begin{enumerate}
\item Define an equilibrium in this economy.

\item Solve for the competitive equilibrium prices and quantities.

\item Now suppose $y=zk+zn.$ Solve for the competitive equilibrium prices
and quantities. Here you will need to be careful about the requirement that $%
\ell \leq 1.$
\end{enumerate}
\end{enumerate}
\end{document}
